% GENERATED FILE. Edit canonical lessons, then run npm run generate:books.
% canonical-source-hash: fnv1a64:f061bebb77ef1225
% canonical-lessons: RU-C194-pelmeni, RU-C194-salat, RU-C194-kapusta, RU-C194-svyokla, RU-C194-gorokh

\chapter{Dumplings, Salad, Cabbage, Beetroot, Peas}
\label{ch:ru-dumplings-salad-cabbage-beetroot-peas}
\hlchaptermodality{194}

\begin{chapteropening}
\textbf{By the end of this chapter:} \emph{I can say dumplings, a salad, cabbage, beetroot and peas.}

Complete the last lesson of chapter 194: I can say dumplings, a salad, cabbage, beetroot and peas.
\end{chapteropening}

\section[\texorpdfstring{\ru{пельмени} (pelméni)}{пельмени (pelméni)}]{\ru{пельмени} (pelméni) — dumplings}

\label{lesson:RU-C194-pelmeni}



\begin{quote}
Before the new one: say the Russian for a button, then the Russian for earrings.
\end{quote}

\subsection*{You'll want to know: \ru{пельмени}}
\textbf{\ru{пельмени}} (\emph{pelméni}) — \textquotedblleft{}dumplings\textquotedblright{}.

\begin{grammarlens}[title={What you've built}]
One more word for this chapter.
\end{grammarlens}

\subsection*{Your turn}
\begin{itemize}
\raggedright
  \item \emph{Say it:} \emph{pelméni}
  \item \emph{Say it:} \emph{pelméni}, once more
  \item \emph{Say it:} say \emph{pelméni}
  \item \emph{Recall:} say the Russian for a button, then the Russian for earrings, then say \emph{pelméni} again
\end{itemize}

\begin{tcolorbox}[breakable,colback=teal!4,colframe=teal!35!black,title={Before you move on}]
What does \ru{пельмени} mean? (\textquotedblleft{}Dumplings\textquotedblright{}.) Say it once more.
\end{tcolorbox}
\section[\texorpdfstring{\ru{салат} (salát)}{салат (salát)}]{\ru{салат} (salát) — a salad}

\label{lesson:RU-C194-salat}



\begin{quote}
Before the new one: say the Russian for earrings, then the Russian for dumplings.
\end{quote}

\subsection*{You'll want to know: \ru{салат}}
\textbf{\ru{салат}} (\emph{salát}) — \textquotedblleft{}a salad\textquotedblright{}.

\begin{grammarlens}[title={What you've built}]
One more word for this chapter.
\end{grammarlens}

\subsection*{Your turn}
\begin{itemize}
\raggedright
  \item \emph{Say it:} \emph{salát}
  \item \emph{Say it:} \emph{salát}, once more
  \item \emph{Say it:} say \emph{salát}
  \item \emph{Recall:} say the Russian for earrings, then the Russian for dumplings, then say \emph{salát} again
\end{itemize}

\begin{tcolorbox}[breakable,colback=teal!4,colframe=teal!35!black,title={Before you move on}]
What does \ru{салат} mean? (\textquotedblleft{}A salad\textquotedblright{}.) Say it once more.
\end{tcolorbox}
\section[\texorpdfstring{\ru{капуста} (kapústa)}{капуста (kapústa)}]{\ru{капуста} (kapústa) — cabbage}

\label{lesson:RU-C194-kapusta}



\begin{quote}
Before the new one: say the Russian for dumplings, then the Russian for a salad.
\end{quote}

\subsection*{You'll want to know: \ru{капуста}}
\textbf{\ru{капуста}} (\emph{kapústa}) — \textquotedblleft{}cabbage\textquotedblright{}.

\begin{grammarlens}[title={What you've built}]
One more word for this chapter.
\end{grammarlens}

\subsection*{Your turn}
\begin{itemize}
\raggedright
  \item \emph{Say it:} \emph{kapústa}
  \item \emph{Say it:} \emph{kapústa}, once more
  \item \emph{Say it:} say \emph{kapústa}
  \item \emph{Recall:} say the Russian for dumplings, then the Russian for a salad, then say \emph{kapústa} again
\end{itemize}

\begin{tcolorbox}[breakable,colback=teal!4,colframe=teal!35!black,title={Before you move on}]
What does \ru{капуста} mean? (\textquotedblleft{}Cabbage\textquotedblright{}.) Say it once more.
\end{tcolorbox}
\section[\texorpdfstring{\ru{свёкла} (svyókla)}{свёкла (svyókla)}]{\ru{свёкла} (svyókla) — beetroot}

\label{lesson:RU-C194-svyokla}



\begin{quote}
Before the new one: say the Russian for a salad, then the Russian for cabbage.
\end{quote}

\subsection*{You'll want to know: \ru{свёкла}}
\textbf{\ru{свёкла}} (\emph{svyókla}) — \textquotedblleft{}beetroot\textquotedblright{}.

\begin{grammarlens}[title={What you've built}]
One more word for this chapter.
\end{grammarlens}

\subsection*{Your turn}
\begin{itemize}
\raggedright
  \item \emph{Say it:} \emph{svyókla}
  \item \emph{Say it:} \emph{svyókla}, once more
  \item \emph{Say it:} say \emph{svyókla}
  \item \emph{Recall:} say the Russian for a salad, then the Russian for cabbage, then say \emph{svyókla} again
\end{itemize}

\begin{tcolorbox}[breakable,colback=teal!4,colframe=teal!35!black,title={Before you move on}]
What does \ru{свёкла} mean? (\textquotedblleft{}Beetroot\textquotedblright{}.) Say it once more.
\end{tcolorbox}
\section[\texorpdfstring{\ru{горох} (gorókh)}{горох (gorókh)}]{\ru{горох} (gorókh) — peas}

\label{lesson:RU-C194-gorokh}



\begin{quote}
Before the new one: say the Russian for cabbage, then the Russian for beetroot.
\end{quote}

\subsection*{You'll want to know: \ru{горох}}
\textbf{\ru{горох}} (\emph{gorókh}) — \textquotedblleft{}peas\textquotedblright{}.

\begin{grammarlens}[title={What you've built}]
One more word for this chapter.
\end{grammarlens}

\subsection*{Your turn}
\begin{itemize}
\raggedright
  \item \emph{Say it:} \emph{gorókh}
  \item \emph{Say it:} \emph{gorókh}, once more
  \item \emph{Say it:} say \emph{gorókh}
  \item \emph{Recall:} say the Russian for cabbage, then the Russian for beetroot, then say \emph{gorókh} again
\end{itemize}

\begin{tcolorbox}[breakable,colback=teal!4,colframe=teal!35!black,title={Before you move on}]
What does \ru{горох} mean? (\textquotedblleft{}Peas\textquotedblright{}.) Say it once more.
\end{tcolorbox}
